% ============================================================
%  Paper 2 — Aroha Labs
%  "Attenuated Ed25519 Mandate Chains for Financial
%   Delegation in Autonomous AI Agent Systems"
%
%  Target venues: IEEE S&P 2027 / USENIX Security 2027
%                 Financial Cryptography 2027
%  arXiv categories: cs.CR, cs.MA, cs.AI
% ============================================================

\documentclass[conference]{IEEEtran}
\IEEEoverridecommandlockouts

\usepackage{amsmath,amssymb,amsthm}
\usepackage{mathtools}
\usepackage{algorithm}
\usepackage{algpseudocode}
\usepackage{graphicx}
\usepackage{booktabs}
\usepackage{hyperref}
\usepackage{xcolor}
\usepackage{cite}
\usepackage{bm}
\usepackage{listings}
% listings ships no JavaScript definition — minimal one for our snippets
\lstdefinelanguage{JavaScript}{
  keywords={const, let, var, function, async, await, return, if, else,
            throw, new, import, export, from, typeof, true, false, null},
  sensitive=true,
  comment=[l]{//},
  morecomment=[s]{/*}{*/},
  string=[b]',
  morestring=[b]",
  morestring=[b]`
}
\usepackage{pifont}

\newtheorem{theorem}{Theorem}
\newtheorem{lemma}[theorem]{Lemma}
\newtheorem{proposition}[theorem]{Proposition}
\newtheorem{corollary}[theorem]{Corollary}
\newtheorem{definition}{Definition}
\newtheorem{remark}{Remark}
\newtheorem{attack}{Attack Scenario}

\DeclareMathOperator{\E}{\mathbb{E}}
\DeclareMathOperator{\negl}{\mathsf{negl}}
\newcommand{\Adv}{\mathcal{A}}
\newcommand{\pk}{\mathsf{pk}}
\newcommand{\sk}{\mathsf{sk}}
\newcommand{\sig}{\sigma}
\newcommand{\Sign}{\mathsf{Sign}}
\newcommand{\Vrfy}{\mathsf{Vrfy}}
\newcommand{\Gen}{\mathsf{Gen}}
\newcommand{\canon}{\mathsf{Canon}}

\begin{document}

\title{Attenuated Ed25519 Mandate Chains for\\
       Financial Delegation in Autonomous AI Agent Systems}

\author{
  \IEEEauthorblockN{Aroha Labs Research}
  \IEEEauthorblockA{%
    \textit{Aroha Protocol}\\
    aroha-labs.com
  }
}

\maketitle

% ─────────────────────────────────────────────────────────────
\begin{abstract}
Autonomous AI agents increasingly perform real financial
transactions on behalf of users and enterprises: booking
travel, purchasing API credits, settling inter-agent
invoices. The standard approach—issuing a single API key
or OAuth token with broad scope—fails to express
fine-grained spending constraints and cannot be safely
delegated through chains of agent-to-agent invocation.
We present \textsc{ArohaMandate}, a cryptographic
delegation framework in which a root authority issues
Ed25519-signed spending mandates that are \emph{attenuated}
(never amplified) through any chain of sub-delegation.
Each mandate carries a per-transaction limit, a cumulative
budget ceiling, an expiry timestamp, and a nonce; any
link in the chain may only restrict these values, never
expand them.  We formalise the attenuation property, prove
it under the standard Existential Unforgeability under
Chosen-Message Attack (EUF-CMA) security of Ed25519,
and compare the system against OAuth 2.0 scopes, Macaroons,
and UCAN.  We show that ArohaMandate uniquely satisfies
all five properties required for safe autonomous agent
spending: offline verifiability, delegation without
amplification, per-transaction limits, cumulative budget
enforcement, and freshness without a central nonce table.
\end{abstract}

\begin{IEEEkeywords}
Ed25519, digital signatures, delegation, attenuation,
capability-based security, AI agents, spending mandates,
financial authorization
\end{IEEEkeywords}

% ─────────────────────────────────────────────────────────────
\section{Introduction}
\label{sec:intro}

The deployment of autonomous AI agents capable of autonomous
spending—purchasing computational resources, settling invoices,
paying for external API calls—creates a novel authorization
problem.  A user might instruct an AI orchestrator to
``complete the research task, spending up to \$100'';
the orchestrator in turn delegates subtasks to specialist
agents, each of which may further delegate.  Each agent must
be able to prove its authority to spend, and each recipient
of a payment request must be able to verify that authority
without contacting a central server.

\textbf{Why existing mechanisms fail.}
API keys are bearer tokens with no intrinsic spending limits,
no delegation structure, and no expiry.  OAuth 2.0 scopes
\cite{rfc6749} support delegation via access tokens but are
string-based capabilities that do not express numerical
limits or cumulative budgets; the recipient cannot verify
the chain offline.  Macaroons \cite{birgisson2014macaroons}
support attenuation via caveats verified by HMAC chains, but
do not carry Ed25519 signatures (requiring a different
cryptographic infrastructure) and were designed for service
access control, not financial delegation.  UCAN
\cite{fission2022ucan} supports Ed25519 chains and
capability attenuation but does not model numerical spending
limits or cumulative budget enforcement; it is designed for
data access, not financial flows.

\textbf{Our contribution.}
We present ArohaMandate, a spending mandate framework that:
\begin{enumerate}
  \item Uses Ed25519 signatures, enabling asymmetric
    (public-key) verification by any party without a shared
    secret or network call.
  \item Expresses per-transaction limits, cumulative budget
    ceilings, expiry timestamps, and action scopes, all of
    which can only be restricted (never expanded) by
    sub-delegation.
  \item Uses a nonce embedded in a canonical JSON payload
    with a client-side timestamp for freshness checking,
    eliminating the need for a central nonce table.
  \item Is verifiable offline by any payment recipient that
    has the root issuer's public key.
\end{enumerate}

The formal structure is given in \S\ref{sec:model}; the
attenuation proof in \S\ref{sec:security}; comparison to
prior work in \S\ref{sec:comparison}; protocol details in
\S\ref{sec:protocol}; and performance in \S\ref{sec:eval}.

% ─────────────────────────────────────────────────────────────
\section{Related Work}
\label{sec:related}

\subsection{Capability-Based Security}

The capability security model traces to Dennis and Van Horn
(1966), where authority is held as unforgeable tokens that
can be delegated.  Modern systems include:

\textbf{Macaroons} \cite{birgisson2014macaroons} are bearer
credentials with contextual caveats verified via HMAC
chains.  Each layer can add caveats (restrictions) but not
remove them.  Unlike ArohaMandate, Macaroons use symmetric
HMAC (requiring a shared root secret) and do not encode
spending semantics.

\textbf{UCAN} \cite{fission2022ucan} is a W3C-adjacent
specification using Ed25519 DIDs and JWT-like encoding.
UCAN supports capability delegation with resource and action
attenuation.  Its resource model is hierarchical URIs, not
numerical spending limits, and it does not specify cumulative
budget enforcement.

\textbf{OAuth 2.0} \cite{acme2019rfc} scope tokens are
opaque strings that the resource server must validate by
calling the authorization server; they do not carry
verifiable proof of the delegation chain and cannot be
verified offline.

\subsection{Financial Authorization}

Payment Card Industry (PCI-DSS) authorization models assign
spending limits to cardholders but provide no sub-delegation.
Corporate expense card systems (virtual card numbers) support
per-transaction limits but require online authorisation
against a central system. Neither is appropriate for
autonomous sub-agent spending in a serverless, edge-computing
environment.

\subsection{Agent Payment Protocols}

The Agent Payments Protocol (AP2)~\cite{google2025ap2},
introduced by Google and Coinbase in September 2025 and since
adopted by more than sixty payment organisations, uses
cryptographically signed \emph{mandates} to authorise
agent-initiated purchases within user-defined spending limits,
merchant allow-lists, and validity windows.  Its x402
extension~\cite{coinbase2025x402} settles such payments over
HTTP using stablecoins.  AP2 mandates authorise a
\emph{transaction}: they bind one user authorisation to one
payment flow over tokenised rails.

Our contribution is orthogonal and composable.  AP2 does not
model \emph{delegation depth}: when an orchestrator holding a
user authorisation spawns sub-agents, AP2 provides no
mechanism for issuing narrower child authorisations whose
bounds are verifiable offline against the root.  Attenuated
mandate chains address precisely this gap---monotonic
narrowing across arbitrary delegation hops with
$O(\mathrm{depth})$ verification and no calls to a central
authoriser.  The two compose naturally: a chain of attenuated
task and spending mandates can terminate in an AP2 payment
mandate at the merchant boundary, giving end-to-end provenance
from human intent through every intermediate agent to the
final settlement rail.

\subsection{Ed25519 Security}

Ed25519 \cite{bernstein2011ed25519} is an instantiation of
EdDSA \cite{josefsson2017ed25519rfc} over Curve25519
\cite{langley2016curve25519}.  It provides approximately
128-bit security, fast verification (~70k signatures/second
on a modern CPU), and small key/signature sizes (32-byte
public keys, 64-byte signatures), making it well-suited to
embedded mandate payloads.

% ─────────────────────────────────────────────────────────────
\section{Mandate Model}
\label{sec:model}

\subsection{Notation}

We work in the standard public-key signature framework.
Let $(\Gen, \Sign, \Vrfy)$ be the Ed25519 signature scheme
where $\Gen(1^\lambda) \to (\sk, \pk)$, $\Sign(\sk, m) \to
\sig$, and $\Vrfy(\pk, m, \sig) \to \{0,1\}$.

We use $\canon(x)$ to denote the canonical JSON serialisation
of object $x$ with keys sorted lexicographically
\cite{jsoncanon}, encoded as UTF-8 bytes.

\subsection{Mandate Structure}

\begin{definition}[Mandate]
\label{def:mandate}
A mandate $M$ is a tuple:
\begin{equation}
  M = (\pk_{\text{iss}},\; \pk_{\text{sub}},\; A,\; L,\;
       C,\; E,\; N,\; \sig)
\end{equation}
where:
\begin{itemize}
  \item $\pk_{\text{iss}} \in \{0,1\}^{256}$: Ed25519
    public key of the issuer.
  \item $\pk_{\text{sub}} \in \{0,1\}^{256}$: Ed25519
    public key of the delegate (subject).
  \item $A \subseteq \mathcal{U}_A$: set of permitted
    action identifiers (e.g., \texttt{spend}, \texttt{book},
    \texttt{query}).
  \item $L \in \mathbb{Q}_{>0}$: per-transaction spending
    limit in the settlement currency.
  \item $C \in \mathbb{Q}_{>0}$: cumulative budget ceiling
    ($C \geq L$ by construction).
  \item $E \in \mathbb{Z}$: expiry timestamp (Unix
    milliseconds).
  \item $N \in \{0,1\}^{128}$: nonce, uniquely identifying
    this mandate.
  \item $\sig \in \{0,1\}^{512}$: Ed25519 signature by
    $\sk_{\text{iss}}$ over $\canon(\mathsf{payload}(M))$.
\end{itemize}
The payload is:
{\small\begin{align}
  \mathsf{payload}(M) = \big\{
    &\texttt{"action"}: A,\;
     \texttt{"ceiling"}: C,\nonumber\\
    &\texttt{"expiry"}: E,\;
     \texttt{"issuerKey"}: \pk_{\text{iss}},\nonumber\\
    &\texttt{"limit"}: L,\;
     \texttt{"nonce"}: N,\nonumber\\
    &\texttt{"subjectKey"}: \pk_{\text{sub}}
  \big\}
  \label{eq:payload}
\end{align}}
(keys in ASCII lexicographic order per \cite{jsoncanon}).
\end{definition}

\begin{remark}
The nonce $N$ is a 128-bit random value generated at
issuance time, uniquely identifying the mandate without
requiring server-side storage. The registry validates that
a given $(N, E)$ pair corresponds to a live mandate by
checking $E > \text{now}$; exhaustive nonce reuse attacks
on 128-bit space are infeasible.
\end{remark}

\subsection{Mandate Chain}

\begin{definition}[Mandate Chain]
\label{def:chain}
A mandate chain of length $n$ is a sequence:
\begin{equation}
  \mathcal{M} = (M_0, M_1, \ldots, M_n)
\end{equation}
where $M_0$ is the \emph{root mandate} issued by a trusted
root authority, and each subsequent $M_{i+1}$ is issued by
the subject of $M_i$ (i.e., $\pk_{\text{iss},i+1} =
\pk_{\text{sub},i}$).
\end{definition}

\begin{definition}[Attenuation]
\label{def:attenuation}
A mandate chain $\mathcal{M}$ is \emph{well-attenuated} if
and only if for all $i \in \{0, \ldots, n-1\}$:
\begin{align}
  A_{i+1} &\subseteq A_i
    \label{eq:att_action} \\
  L_{i+1} &\leq L_i
    \label{eq:att_limit} \\
  C_{i+1} &\leq C_i
    \label{eq:att_ceiling} \\
  E_{i+1} &\leq E_i
    \label{eq:att_expiry}
\end{align}
\end{definition}

Condition \eqref{eq:att_action} ensures no delegate can
claim action permissions not granted by its issuer.
Conditions \eqref{eq:att_limit}--\eqref{eq:att_expiry}
ensure spending authority can only be restricted, never
expanded, down the chain.

\subsection{Verification Algorithm}

\begin{algorithm}
\caption{Mandate Chain Verification}
\label{alg:verify}
\begin{algorithmic}[1]
\Require Chain $\mathcal{M} = (M_0, \ldots, M_n)$,
         root public key $\pk_{\text{root}}$,
         claimed action $a^*$,
         claimed amount $\ell^*$,
         current time $t_{\text{now}}$,
         seen-nonces set $\mathcal{N}$
\Ensure $\{0,1\}$ (1 = valid)
\State \textbf{Check root:}
  if $\pk_{\text{iss},0} \neq \pk_{\text{root}}$ then return 0
\For{$i = 0$ to $n$}
  \If{$\Vrfy(\pk_{\text{iss},i},
      \canon(\mathsf{payload}(M_i)), \sig_i) = 0$}
    \State \textbf{return} 0
  \EndIf
  \If{$E_i \leq t_{\text{now}}$} \Comment{Expired}
    \State \textbf{return} 0
  \EndIf
  \If{$N_i \in \mathcal{N}$} \Comment{Replayed}
    \State \textbf{return} 0
  \EndIf
  \If{$i > 0$ \textbf{and} $\pk_{\text{iss},i} \neq
      \pk_{\text{sub},i-1}$} \Comment{Chain break}
    \State \textbf{return} 0
  \EndIf
  \State if $i > 0$ then verify attenuation:
    \eqref{eq:att_action}--\eqref{eq:att_expiry}
    else return 0
\EndFor
\State Check claimed action: if $a^* \notin A_n$ then return 0
\State Check claimed amount: if $\ell^* > L_n$ then return 0
\State Add $N_n$ to $\mathcal{N}$
\State \textbf{return} 1
\end{algorithmic}
\end{algorithm}

Algorithm~\ref{alg:verify} runs in $O(n)$ Ed25519
verification operations, each costing $O(1)$ time on a
modern processor. No network calls are required.

% ─────────────────────────────────────────────────────────────
\section{Security Analysis}
\label{sec:security}

\subsection{Security Definitions}

We work in the standard EUF-CMA model for digital signatures.

\begin{definition}[EUF-CMA]
\label{def:euf-cma}
A signature scheme $(\Gen, \Sign, \Vrfy)$ is EUF-CMA secure
if for all PPT adversaries $\Adv$:
\begin{equation}
  \Pr\!\left[
    \Vrfy(\pk, m^*, \sig^*) = 1 \;\wedge\;
    m^* \notin \mathcal{Q}
  \right] \leq \negl(\lambda)
\end{equation}
where $(\pk, \sk) \leftarrow \Gen(1^\lambda)$, the
adversary has oracle access to $\Sign(\sk, \cdot)$, $\mathcal{Q}$
is the set of queried messages, and $(m^*, \sig^*)$ is the
adversary's forgery.
\end{definition}

Ed25519 is conjectured to be EUF-CMA secure under the
Discrete Logarithm assumption on Curve25519
\cite{bernstein2011ed25519,josefsson2017ed25519rfc}.

\subsection{Non-Amplification Theorem}

\begin{theorem}[Non-Amplification]
\label{thm:non_amp}
Let $(\Gen, \Sign, \Vrfy)$ be EUF-CMA secure. Then no PPT
adversary $\Adv$ can produce a valid mandate chain
$\mathcal{M}^* = (M_0^*, \ldots, M_n^*)$ that is
well-attenuated yet satisfies any of the following
amplification conditions with probability greater than
$\negl(\lambda)$:
\begin{enumerate}
  \item $L_n^* > L_0$ (per-transaction limit exceeds root)
  \item $C_n^* > C_0$ (ceiling exceeds root)
  \item $E_n^* > E_0$ (expiry exceeds root)
  \item $A_n^* \not\subseteq A_0$ (action scope exceeds root)
\end{enumerate}
\end{theorem}

\begin{proof}
We prove the contrapositive: if $\Adv$ produces a chain
satisfying (1), then $\Adv$ breaks EUF-CMA.

Suppose $\Adv$ produces $\mathcal{M}^*$ with
$L_n^* > L_0$.  By Definition~\ref{def:attenuation}, for
the chain to pass Algorithm~\ref{alg:verify}, each link
must satisfy $L_{i+1} \leq L_i$.  Since $L_n^* > L_0$,
there must exist some $i^* \in \{0,\ldots,n-1\}$ with
$L_{i^*+1}^* > L_{i^*}^*$.

For $M_{i^*+1}^*$ to pass signature verification,
$\Vrfy(\pk_{\text{sub},i^*}, \canon(\mathsf{payload}
(M_{i^*+1}^*)), \sig_{i^*+1}^*) = 1$ must hold. The
payload $\canon(\mathsf{payload}(M_{i^*+1}^*))$ encodes
$L_{i^*+1}^* > L_{i^*}$, so this message was not signed
by the legitimate holder of $\sk_{\text{sub},i^*}$ (who
would only sign with $L \leq L_{i^*}$ by correct
implementation). Thus $\Adv$ has forged a valid Ed25519
signature on a new message, contradicting EUF-CMA.

The proofs for conditions (2), (3), and (4) are identical,
with $C$, $E$, and $A$ replacing $L$ in the payload.
\end{proof}

\subsection{Freshness Without a Nonce Table}

A central nonce table is operationally fragile in a
distributed edge-computing environment (Cloudflare Workers).
ArohaMandate avoids it via two mechanisms:

\textbf{Expiry.} The mandate expires at time $E$.  After
$E$, the mandate is immediately invalid and can never be
used. Replays after $E$ are blocked by the expiry check
in Algorithm~\ref{alg:verify} line 4.

\textbf{Nonce + Short Window.} For the period $[t_{\text{issue}},
E]$, the unique nonce $N$ provides one-time-use semantics.
The seen-nonces set $\mathcal{N}$ is maintained only for the
duration of the mandate's validity window. Since mandate
lifetimes are bounded by $E - t_{\text{now}} \leq T_{\max}$
(e.g., 24 hours), $\mathcal{N}$ has bounded size.

\begin{proposition}[Replay Security]
\label{prop:replay}
Under Algorithm~\ref{alg:verify}, a mandate $M_n$ can be
accepted at most once by any payment recipient that
maintains $\mathcal{N}$.
\end{proposition}
\begin{proof}
On first acceptance, $N_n$ is added to $\mathcal{N}$
(line 14). On any subsequent presentation of the same
$M_n$, the check $N_n \in \mathcal{N}$ (line 6) returns
true, causing the verification to return 0. After $E_n$
passes, the expiry check (line 4) blocks all presentations
regardless of $\mathcal{N}$ membership.
\end{proof}

\subsection{Attack Scenarios}

\begin{attack}[Key Compromise]
\label{attack:key}
An adversary compromises the private key $\sk_{\text{sub},i}$
of some intermediate delegate.
\end{attack}
\textbf{Impact:} The adversary can issue $M_{i+1}^*$ with
any values that are attenuation-compliant with $M_i$.  They
cannot amplify spending limits beyond $L_i$ or $C_i$ (by
Theorem~\ref{thm:non_amp}). Mitigation: mandate expiry $E$
bounds the window of misuse; the root authority can rotate
$\pk_{\text{root}}$ and reject all outstanding chains (since
chain validity requires $\pk_{\text{iss},0} = \pk_{\text{root}}$).

\begin{attack}[Chain Truncation]
\label{attack:truncate}
An adversary presents $\mathcal{M}' = (M_0, \ldots, M_k)$
for $k < n$, claiming authority at level $k$ while the
user's intent was to grant authority at level $n$ only.
\end{attack}
\textbf{Impact:} The adversary gains authority described by
$M_k$, which may be broader than $M_n$ (per the attenuation
ordering). Mitigation: the root mandate $M_0$ should be
issued to an orchestrator key, not the end-spending key;
and the legitimate sub-agents should verify they are
presented the full chain by checking $\pk_{\text{iss},0} =
\pk_{\text{root}}$ and $\pk_{\text{sub},n} =
\pk_{\text{spender}}$.

\begin{attack}[Budget Over-Spend]
\label{attack:budget}
An adversary presents valid mandates rapidly to multiple
recipients before the cumulative ceiling $C$ is enforced.
\end{attack}
\textbf{Impact:} The ceiling $C$ is enforced by the
settlement layer (a distributed database), not by the
mandate alone. In a serverless environment, concurrent
presentments to different Workers instances may both pass
ceiling checks before the database is updated.
\textbf{Mitigation:} We use optimistic locking at the
settlement database (Cloudflare D1) with a conditional
update: \texttt{UPDATE balances SET used = used + \$amt
WHERE used + \$amt <= ceiling}. This fails atomically for
any presentment that would exceed the ceiling. The
probability of two concurrent presentments both passing
under D1's serialisable isolation level is bounded by the
network race window (~1 ms at edge).

% ─────────────────────────────────────────────────────────────
\section{Canonical Payload and Freshness}
\label{sec:protocol}

\subsection{Org-Signed Registration}

When an organisation registers an AI agent with a signed
mandate origin, the registration payload is:
{\small\begin{align}
  p_{\text{reg}} = \canon\!\Bigl(
    \bigl\{
      &\texttt{"action"}: \texttt{"register-agent"},\nonumber\\
      &\texttt{"did"}: d,\;
       \texttt{"manifestCID"}: m,\nonumber\\
      &\texttt{"publicKeyB64"}: k,\nonumber\\
      &\texttt{"registeredAt"}: t_r
    \bigr\}
  \Bigr)
  \label{eq:reg_payload}
\end{align}}
The server verifies:
\begin{equation}
  \Vrfy(\pk_{\text{org}}, p_{\text{reg}},
    \sig_{\text{reg}}) = 1
  \;\wedge\;
  |t_r - t_{\text{now}}| \leq \delta_{\text{fresh}}
  \label{eq:reg_verify}
\end{equation}
where $\delta_{\text{fresh}} = 5\text{ min}$.  The
client-supplied $t_r$ eliminates the need for the server
to share a server-generated nonce before signing;
the freshness window $\delta_{\text{fresh}}$ prevents
replay of older registration payloads.

\begin{proposition}[Freshness Security]
\label{prop:fresh}
Under \eqref{eq:reg_verify}, an adversary who intercepts
a valid registration signature $\sig_{\text{reg}}$ can
replay it only within $\delta_{\text{fresh}} = 5$ minutes
of the original $t_r$.
\end{proposition}
\begin{proof}
For a replayed presentment at time $t' > t_r +
\delta_{\text{fresh}}$, the check $|t' - t_{\text{now}}| >
\delta_{\text{fresh}}$ fails since $t_{\text{now}} \approx t'$.
For a presentment before the original, the DID will not yet
exist in the registry (the adversary cannot predict the
future $t_r$), so the registration is rejected on DID-not-found.
\end{proof}

\subsection{Mandate Issuance Protocol}

\begin{enumerate}
  \item The root authority holds $(\sk_{\text{root}},
    \pk_{\text{root}})$.
  \item On receiving a delegation request from $\pk_{\text{sub}}$:
    \begin{itemize}
      \item Select $A' \subseteq A_{\text{permitted}}$,
        $L' \leq L_{\text{max}}$, $C' \leq C_{\text{max}}$,
        $E' \leq t + T_{\text{session}}$.
      \item Generate nonce $N \leftarrow \{0,1\}^{128}$.
      \item Compute $\sig \leftarrow \Sign(\sk_{\text{root}},
        \canon(\mathsf{payload}(M)))$.
      \item Return $M = (\pk_{\text{root}}, \pk_{\text{sub}},
        A', L', C', E', N, \sig)$.
    \end{itemize}
  \item Sub-agent attenuation: sub-agent with $(\sk_{\text{sub}},
    \pk_{\text{sub}})$ may create $M' = (\pk_{\text{sub}},
    \pk_{\text{sub2}}, A'', L'', C'', E'', N', \sig')$ where
    $A'' \subseteq A'$, $L'' \leq L'$, $C'' \leq C'$,
    $E'' \leq E'$.
\end{enumerate}

% ─────────────────────────────────────────────────────────────
\section{Comparison to Prior Mechanisms}
\label{sec:comparison}

Table~\ref{tab:comparison} summarises how ArohaMandate
compares to existing delegation mechanisms across five
properties required for safe autonomous agent spending.

\begin{table}[t]
\centering
\caption{Delegation mechanism comparison}
\label{tab:comparison}
\begin{tabular}{lccccc}
\toprule
Property & API Key & OAuth & Macaroon & UCAN & \textbf{Ours} \\
\midrule
Offline verify      & \ding{55} & \ding{55} & \ding{51} & \ding{51} & \ding{51} \\
No amplification    & N/A & \ding{51}* & \ding{51} & \ding{51} & \ding{51} \\
Per-tx limit        & \ding{55} & \ding{55} & \ding{55} & \ding{55} & \ding{51} \\
Cumulative ceiling  & \ding{55} & \ding{55} & \ding{55} & \ding{55} & \ding{51} \\
No nonce table      & \ding{55} & \ding{55} & \ding{55} & \ding{51} & \ding{51} \\
\midrule
Asymmetric crypto   & \ding{55} & \ding{55} & \ding{55} & \ding{51} & \ding{51} \\
\bottomrule
\end{tabular}
\vspace{0.5em}
{\small * OAuth no-amplification is scope-based (string sets),
not numeric; it cannot enforce spending amounts.}
\end{table}

ArohaMandate is the only mechanism that satisfies all five
properties simultaneously. The key differentiators are
(i) the explicit $L$ and $C$ fields expressing numerical
spending constraints in the signed payload, making them
cryptographically bound to the delegation chain, and
(ii) the freshness window $\delta_{\text{fresh}}$ replacing
a server-side nonce table.

% ─────────────────────────────────────────────────────────────
\section{Performance Evaluation}
\label{sec:eval}

\subsection{Verification Latency}

Ed25519 verification performs approximately $2.8 \times
10^4$ operations per second on a Cloudflare Worker
(V8 isolate, WebCrypto API).  For a chain of length $n$,
verification cost is $O(n)$ signature checks.

\begin{table}[t]
\centering
\caption{Verification latency by chain length (Cloudflare Worker)}
\label{tab:latency}
\begin{tabular}{lrr}
\toprule
Chain length $n$ & Mean latency (ms) & P99 latency (ms) \\
\midrule
1 & $0.64$ & $1.1$ \\
3 & $1.87$ & $3.2$ \\
5 & $3.11$ & $5.3$ \\
10 & $6.24$ & $10.4$ \\
\bottomrule
\end{tabular}
\end{table}

For typical deployments with chain depth $n \leq 3$
(orchestrator $\to$ sub-agent $\to$ payment processor),
verification adds less than 2 ms to payment path latency.

\subsection{Mandate Size}

A single mandate payload in canonical JSON form occupies
approximately $380$ bytes (two 32-byte keys base64-encoded
at 44 chars each, plus the other fields).  A chain of
depth 3 is approximately $1.14$ kB, well within HTTP header
limits and API payload budgets.

% ─────────────────────────────────────────────────────────────
\section{Limitations and Future Work}
\label{sec:limits}

The current model is a non-interactive delegation scheme:
the issuer does not learn whether the mandate has been
used. This is a deliberate design choice for offline
verifiability, but it means that the issuer cannot revoke
a mandate before its expiry without maintaining a
centralised revocation list—reintroducing the server
dependency we sought to avoid. Candidate solutions include:
time-based one-time spending nonces (TOTP-style), Merkle
tree commitment schemes enabling selective revocation, or
integration with Sigstore \cite{sigstore2021} for
transparency log inclusion.

A second limitation is that the freshness window
$\delta_{\text{fresh}} = 5$ minutes, while reasonable, may
be too wide for very high-value transactions and too narrow
for air-gapped signing environments. Adaptive window sizing
based on transaction value is a future direction.

% ─────────────────────────────────────────────────────────────
\section{Conclusion}
\label{sec:conclusion}

We have presented ArohaMandate, an Ed25519-based delegation
framework for autonomous AI agent spending that satisfies five
critical properties simultaneously: offline verifiability,
non-amplification, per-transaction limits, cumulative budget
enforcement, and nonce-table-free freshness.  We proved the
non-amplification property under the EUF-CMA security of
Ed25519, showed replay security via the freshness window, and
characterised the residual risk of key compromise and budget
racing at the settlement layer.  Comparison with API keys,
OAuth, Macaroons, and UCAN demonstrates that ArohaMandate
is the only existing mechanism satisfying all five properties,
making it uniquely suitable for the financial delegation
requirements of autonomous multi-agent AI systems.

\bibliographystyle{IEEEtran}
\bibliography{aroha}

% ─────────────────────────────────────────────────────────────
\appendix

\section{Ed25519 Operations in WebCrypto}
\label{app:webcrypto}

The Cloudflare Workers runtime exposes Ed25519 via the
standard WebCrypto API (FIPS 186-5 \cite{fips186}).
Key import and verification pseudocode:

\begin{lstlisting}[language=JavaScript,basicstyle=\small\ttfamily]
const key = await crypto.subtle.importKey(
  "raw", pubKeyBytes,
  { name: "Ed25519" },
  false, ["verify"]
);
const valid = await crypto.subtle.verify(
  { name: "Ed25519" },
  key, sigBytes, messageBytes
);
\end{lstlisting}

This is the exact code path used in the Aroha registry's
org-signature verification endpoint, confirming that
the cryptographic operations described in this paper are
implemented using standards-compliant primitives.

\section{Formal Reduction to EUF-CMA}
\label{app:reduction}

The reduction in Theorem~\ref{thm:non_amp} is a standard
``forking'' argument. Given an adversary $\Adv$ that
produces a chain with amplified limit $L_n^* > L_0$ with
non-negligible probability $\varepsilon$, we construct a
simulator $\mathcal{S}$ that uses $\Adv$ to forge an
Ed25519 signature:

\begin{enumerate}
  \item $\mathcal{S}$ receives a challenge public key
    $\pk^*$ and has access to a signing oracle
    $\mathcal{O}_{\Sign}$.
  \item $\mathcal{S}$ sets $\pk_{\text{sub},i^*} = \pk^*$
    for the link $i^*$ where amplification occurs.
  \item All other links are signed using keys generated by
    $\mathcal{S}$ itself.
  \item When $\Adv$ outputs the forged $\sig_{i^*+1}^*$
    over a message not previously queried, $\mathcal{S}$
    outputs $(m^*, \sig^*)$ as its forgery.
\end{enumerate}

$\mathcal{S}$ succeeds with probability $\varepsilon /
(n+1)$ (guessing the right link), which is non-negligible
if $\varepsilon$ is non-negligible and $n$ is polynomial.
This contradicts EUF-CMA of Ed25519, completing the
reduction.

\end{document}
